\subsection{Colorado}

Colorado's Amendment Y (2018) added Article V, Section 44 to the Colorado
Constitution, establishing an Independent Congressional Redistricting Commission.
The amendment requires the commission to, among other criteria, make congressional
districts ``as compact as practicable.''
The commission's 2021 procedural guidelines operationalised compactness using
a Schwartzberg-equivalent metric, specifying that districts should minimise
the ratio of perimeter to the square root of area~\citep{co_irc2021}.
This is precisely the Schwartzberg score (up to a constant factor of $2\sqrt{\pi}$).

\paragraph{Practical impact.}
Colorado redistricting filings that cite bisect results should report S values
(or equivalently, PP values with a note converting to S).
For reference: the Colorado 2021 congressional map (8 districts) achieved
mean PP $\approx 0.31$--$0.38$ according to Princeton Gerrymandering Project analyses,
corresponding to mean S $\approx 1.63$--$1.79$.

\subsection{Iowa}

Iowa Code \S 42.4 requires that congressional districts be ``as compact as possible.''
The Iowa Legislative Services Agency (LSA) has operationalised this requirement
using a perimeter-to-area ratio that is equivalent to the Schwartzberg score.
Iowa's plans are evaluated using the LSA's compactness formula, which the LSA
has internally standardised to the S metric.

\subsection{Oregon}

Oregon's redistricting statute (ORS 188.010) requires legislative districts to
be ``as compact as possible.''
The Oregon Secretary of State's office has used perimeter-based compactness
measures in its redistricting analysis that are equivalent to S.

\subsection{Utah}

Utah Code Ann.\ \S 20A-19-102 requires redistricting plans to consider compactness.
The Utah Independent Redistricting Commission (2021) used a perimeter-squared
over area metric in its analysis, which is equivalent to the Schwartzberg
score under the transformation $S^2 = P^2/(4\pi A) = 1/\text{PP}$.

\subsection{Summary}

\begin{table}[ht]
\centering
\caption{States with Schwartzberg-equivalent statutory or operational compactness requirements.}
\label{tab:statutory}
\begin{tabular}{lll}
\toprule
State & Statutory reference & S-equivalent operationalisation \\
\midrule
Colorado & Colo.\ Const.\ Art.\ V, \S 44 & Commission guidelines (2021) \\
Iowa     & Iowa Code \S 42.4              & LSA internal standard \\
Oregon   & ORS 188.010                    & Secretary of State practice \\
Utah     & Utah Code \S 20A-19-102        & UIRC methodology (2021) \\
\bottomrule
\end{tabular}
\end{table}
